\chapter{Transporte, selección y prolongación en el núcleo topológico de fase}
\label{sec:arquitectura-tpk}

La sigla TPK designa el \emph{Topological Phase Kernel}, o n\'ucleo
topol\'ogico de fase.\footnote{La denominaci\'on conserva asimismo la
dedicatoria autoral a Tan Peng Kian.} No designa un \`unico algoritmo de 108
pasos, una tabla de resultados ni el cat\'alogo de 468 emisiones. Es el
sistema de operadores que hace convivir, sobre la categor\'ia de caminos de
APP, las lecturas aditiva y multiplicativa sin identificar sus genealog\'ias.

La palabra \emph{topol\'ogico} tiene aqu\'i un contenido preciso. El soporte
de APP es el grafo de Cayley del grupo cociente
\[
 X_9=(\mathbb Z/9\mathbb Z)^2
\]
con sus aristas cardinales. Una ruta cerrada admite una clase de enrollamiento
despu\'es de elevarla al recubridor \(\mathbb Z^2\). TPK transporta esa clase,
la orientaci\'on de la ruta y la holonom\'ia de fase. Dos caminos que terminan
en la misma celda pueden, por tanto, representar estados distintos. El n\'ucleo
de fase es precisamente la informaci\'on que permanece invisible al publicar
solamente la celda final o su residuo.

La base dinámica es la categoría libre \(\mathsf P_{\rm APP}\) del grafo
cardinal de APP\@. Puesto que las dos evaluaciones aritméticas no inducen la
misma partición, el transporte vive sobre la categoría bivariada
\[
 \mathsf P_{\rm APP}^{(2)}
 =\mathsf P_{\rm APP}\times\mathsf P_{\rm APP}.
\]
La primera componente se evalúa mediante la hoja aditiva \(\Sigma\) y la
segunda mediante la hoja multiplicativa \(\Pi\). Esta separación es
constitutiva, pero no duplica APP: registra dos lecturas ordenadas de una
misma matriz de caminos. La componente multiplicativa es la tabla principal;
la aditiva conserva la ley de acumulación que genera sus filas y permite
medir qué información elimina una proyección visible.

\ifdefined\HMTTPKRepeatAPP
\section{Carta canónica y arquitectura suma--producto}

El soporte se escribe en coordenadas de cursor
\[
 I_9=\{0,\ldots,8\},\qquad X_9=I_9^2,
\]
pero sus etiquetas aritm\'eticas pertenecen a \(\{1,\ldots,9\}\). La carta
que relaciona ambos niveles es
\[
 \chi:I_9\longrightarrow\{1,\ldots,9\},\qquad \chi(i)=i+1.
\]
Sea
\[
 \rho_9(n)=1+((n-1)\bmod9),\qquad
 q_9(n)=\frac{n-\rho_9(n)}9.
\]
En particular, \(9\) es el representante visible de la clase nula. Los dos
observables de un mismo v\'ertice son
\[
 \Sigma(i,j)=\rho_9\bigl(\chi(i)+\chi(j)\bigr),\qquad
 \Pi(i,j)=\rho_9\bigl(\chi(i)\chi(j)\bigr).
\]
Esta convenci\'on impide mezclar la coordenada de cursor cero con la etiqueta
visible nueve. APP es la tupla
\[
 \mathcal A_9=
 \bigl(X_9,\operatorname{Cay}(X_9),\Sigma,\Pi,
       \operatorname{Path}(X_9)\bigr),
\]
y no la yuxtaposici\'on de dos tablas: cada semilla \((i,j)\) porta a la vez
las dos evaluaciones y cada ruta conserva el orden en que las consulta.

La elevaci\'on que retiene residuo y cociente es
\[
 \mathcal A_9^\#=(R^+,Q^+;R^\times,Q^\times),
\]
con, para \(a=\chi(i)\) y \(b=\chi(j)\),
\[
 a+b=R^+_{ij}+9Q^+_{ij},\qquad
 ab=R^\times_{ij}+9Q^\times_{ij},
\]
\[
 R^+_{ij}=\Sigma(i,j),\qquad R^\times_{ij}=\Pi(i,j).
\]
Los totales exactos son
\[
 \sum R^+=405,\quad \sum R^\times=459,\quad
 \sum Q^+=45,\quad \sum Q^\times=174.
\]
Por tanto,
\[
 2025-810=1215=(459-405)+9(174-45)=54+9\cdot129.
\]
La igualdad no crea TPK por simple conteo: identifica el dato que el
transporte debe conservar. El residuo visible difiere en \(54\) y la memoria
de cociente en \(129\).

La composici\'on de la hoja aditiva posee el cociclo de acarreo
\[
 \kappa_9(a,b)=
 \frac{\rho_9(a)+\rho_9(b)-\rho_9(a+b)}9,
\]
que satisface
\[
 \kappa_9(a,b)+\kappa_9(a+b,c)
 =\kappa_9(b,c)+\kappa_9(a,b+c).
\]
El sector multiplicativo posee adem\'as el radical
\[
 \mathfrak m=(3)=\{3,6,9\},\qquad \mathfrak m^2=0
 \quad\text{en }\mathbb Z/9\mathbb Z.
\]
Las seis filas unitarias suman \(45\), las filas \(3,6\) suman \(54\) y
la fila \(9\) suma \(81\); as\'i
\[
 459=6\cdot45+2\cdot54+81,
\]
y todo el exceso \(459-405=54\) se localiza en el sector no unitario. Esta
localizaci\'on, el cociclo de acarreo y la diferencia de cocientes son tres
datos distintos de la misma arquitectura suma--producto.
\fi

\section{Reducciones residuales y representación decimal}

La arquitectura usa tres reducciones, cada una con un cometido distinto:
\[
 \mathbb Z\xrightarrow{(\rho_9,q_9)}
 \{1,\ldots,9\}\times\mathbb Z,
 \qquad
 \rho_3^{\rm or}:\mathbb Z/3\mathbb Z\longrightarrow
 \{-1,0,+1\},
\]
\[
 (d_1,d_2,\ldots)\in\{0,\ldots,999\}^{\mathbb N}
 \xrightarrow{\iota_{1000}}
 \sum_{m\geq1}d_m1000^{-m}.
\]
La orientaci\'on ternaria es
\[
 \rho_3^{\rm or}([0])=0,\qquad
 \rho_3^{\rm or}([1])=+1,\qquad
 \rho_3^{\rm or}([2])=-1.
\]
La primera separa residuo visible y cociente non\'adico; la segunda asigna el
tipo orientado de la transici\'on; la tercera agrupa la lectura arquimediana
en tr\'iadas decimales, cuyo registro conserva el acarreo. Ninguna puede
sustituir a otra.
En particular, la reducci\'on m\'odulo tres no reconstruye el cociente m\'odulo
nueve, y una tr\'iada decimal no es una palabra de seis trits.
La suma acumula desplazamientos; el producto pliega clases y contrae las
filas no invertibles. La necesidad de TPK nace de que las dos lecturas
comparten soporte, pero no comparten partici\'on, periodo ni cociente.

\Needspace{36\baselineskip}
\section{Correspondencia entre símbolos del TPK, dominios y composiciones}
\label{sec:diccionario-interno-tpk}

Los nombres históricos se fijan aquí por objeto y no por semejanza
numérica. Esta tabla forma parte de la definición pública de TPK.

\begin{longtable}{@{}>{\raggedright\arraybackslash}p{.19\textwidth}
>{\raggedright\arraybackslash}p{.30\textwidth}
>{\raggedright\arraybackslash}p{.43\textwidth}@{}}
\toprule
Nombre&Objeto&Función dentro del núcleo\\\midrule
\endhead
coordenada nonádica&
\(\rho_9(n)\) junto con \(q_9(n)\)&
Separa el representante visible módulo \(9\) del cociente que conserva la
memoria del cruce de frontera. No es una raíz digital.\\
selector temporal ternario&
\(M_{\rm ph}(t)=\varepsilon(t)\in\{+1,-1,0\}\)&
En \(1,4,7\) activa la acumulación aditiva; en \(2,5,8\) activa la
multiplicativa; en \(3,6,9\) mantiene ambos cursores y registra umbral.
No designa las direcciones espaciales.\\
transporte cardinal&
\(T_N,T_E,T_S,T_O\)&
Desplaza la posición sobre el toro APP y conserva la palabra de ruta y su
enrollamiento. Es independiente del selector temporal.\\
lector ternario&
\(\rho_3^{\rm or}\)&
Convierte las tres clases residuales en orientación
\(0,+1,-1\). No reconstruye por sí solo el cociente nonádico.\\
lector en base mil&
\(\iota_{1000}\)&
Agrupa la publicación arquimediana en tríadas decimales y transporta sus
acarreos. Es un lector posicional, no una congruencia ``módulo \(1000\)''.\\
retorno de veintisiete&
\(\mathcal K_{27}=F_{\rm obs}^{27}\)&
Compone veintisiete actualizaciones posicionales y orientadas. Su cuarto
iterado restaura la fase observable, mientras el registro completo conserva
la memoria acumulada.\\
registro dodecafásico&
\(\Lambda\in\mathcal R_{12}\)&
Conserva seis posiciones, cinco diferencias y un total. Es el soporte del
transporte reversible \(1+5+6=12\).\\
lectura bidireccional&
\(q_+,q_-\) y sus pesos conjugados&
Conserva simultáneamente las orientaciones de ida y retorno. En Mecánica
Dimensional, una firma \((n_A,s_C)\) produce primero
\(W_\pm=6n_A\pm s_C\); sólo después aparece
\(s_C/6\) en el carácter angular.\\
\bottomrule
\end{longtable}

El selector temporal, el transporte cardinal y la orientación ternaria son,
por tanto, tres mapas distintos. Su separación conserva la diferencia entre
\emph{cuándo} se activa una lectura, \emph{hacia dónde} avanza la ruta y
\emph{con qué orientación} se registra el paso. TPK contiene los tres y
conserva sus relaciones.

\ifdefined\HMTTPKWordArithmetic
\section{De la célula local al monoide de palabras nonádicas}

La arquitectura suma--producto no termina en la tabla local. Para probar que
el cociente transportado genera aritm\'etica global, hay que construir la
operaci\'on sobre palabras sin introducir como definici\'on la multiplicaci\'on
entera que se pretende recuperar. Sea
\[
 \mathcal D_9=\{1,\ldots,9\},\qquad
 \mathcal W_9=\{\varnothing\}\sqcup
 \bigsqcup_{\ell\geq1}\mathcal D_9^\ell .
\]
Las palabras se escriben desde la posici\'on menos significativa y se eval\'uan
por
\[
 \operatorname{ev}_9(u_0,\ldots,u_{\ell-1})
 =\sum_{j=0}^{\ell-1}u_j9^j,\qquad
 \operatorname{ev}_9(\varnothing)=0.
\]
El s\'imbolo \(9\) conserva el cruce de frontera; no se sustituye por un cero.

\begin{theorem}[Forma normal non\'adica]
La evaluaci\'on \(\operatorname{ev}_9\) es una biyecci\'on entre
\(\mathcal W_9\) y \(\mathbb N_0\).
\end{theorem}
\begin{proof}
Las palabras de longitud \(\ell\) cubren exactamente el intervalo
\[
 \frac{9^\ell-1}{8}
 \leq\operatorname{ev}_9(u)
 \leq9\frac{9^\ell-1}{8}.
\]
El m\'inimo del intervalo siguiente es una unidad mayor que el m\'aximo del
anterior. Para \(n>0\), la divisi\'on de \(n-1\) por nueve determina de manera
\'unica
\[
 d=1+((n-1)\bmod9)\in\mathcal D_9
\]
y un cociente estrictamente menor, al cual se aplica recursivamente la misma
regla. El proceso termina y es reversible.
\end{proof}
Por ejemplo,
\[
 9\leftrightarrow(9),\quad
 10\leftrightarrow(1,1),\quad
 81\leftrightarrow(9,8),\quad
 1000\leftrightarrow(1,3,3,1).
\]
La \`ultima identidad exhibe en una sola forma normal las dos particiones
fundamentales
\[
 1000=729+271=729+243+27+1;
\]
su interpretaci\'on por canales se introduce despu\'es, no en la definici\'on
de la numeraci\'on.

\section{Suma, producto de Horner y conservación del cociente}

La suma nativa \(u\boxplus_\#v\) recorre las palabras por posiciones. En cada
posici\'on aplica la c\'elula aditiva a los dos d\'igitos presentes y, si existe
un acarreo entrante, aplica una segunda c\'elula al residuo provisional. El
residuo se publica y la suma de cocientes pasa a la posici\'on siguiente. El
invariante parcial es
\[
 \sum_{j<k}(u_j+v_j)9^j
 =\sum_{j<k}w_j9^j+c_k9^k,
\]
de donde se obtiene
\[
 \operatorname{ev}_9(u\boxplus_\#v)
 =\operatorname{ev}_9(u)+\operatorname{ev}_9(v).
\]

Para \(d\in\mathcal D_9\), el transductor escalar
\(d\odot_\#u\) usa primero
\[
 u_jd=9q_j+r_j,\qquad
 (r_j,q_j)=\bigl(R^\times(u_j,d),Q^\times(u_j,d)\bigr),
\]
y normaliza la colisi\'on entre \(r_j\) y el cociente entrante mediante la
c\'elula aditiva. As\'i,
\[
 \operatorname{ev}_9(d\odot_\#u)
 =d\operatorname{ev}_9(u).
\]
Si \(v=(v_0,\ldots,v_{m-1})\), se define, de derecha a izquierda,
\[
 a_m=\varnothing,\qquad
 a_j=(9\odot_\#a_{j+1})\boxplus_\#(v_j\odot_\#u),
\]
y
\[
 u\boxtimes_\#v:=a_0.
\]
Esta es una construcci\'on de Horner realizada s\'olo con las dos c\'elulas
locales y sus cocientes.

Esta descripción fija la dependencia operatoria, pero no constituye una
segunda residencia del resultado. El producto global, su prueba de Horner y
la identidad
\(operatorname{ev}_9(u\boxtimes_\#v)
=\operatorname{ev}_9(u)\operatorname{ev}_9(v)\)
se establecen una sola vez en el
\cref{int:eq:ev-multiplicative}, dentro del capítulo especializado de producto
nativo.
El ejemplo de frontera m\'as corto es
\[
 (9)\boxtimes_\#(9)=(9,8),\qquad 9+8\cdot9=81.
\]
Conservar s\'olo \(R^\times\) destruir\'ia este teorema: \(2\cdot5=10\)
quedar\'ia reducido al residuo visible \(1\). El cociente no es una anotaci\'on
auxiliar, sino la memoria necesaria para que el producto sobreviva al paso por
la frontera.

\section{Coproducto e irreducibilidad}

La apertura multiplicativa se define internamente por
\[
 \Delta_\#e_w
 =\sum_{u\boxtimes_\#v=w}e_u\otimes e_v.
\]
Su versi\'on reducida elimina los factores \((1)\). Si
\(\widetilde\Delta_\#\) denota esa versi\'on y
\[
 A_\#=
 \bigl(\overline{\widetilde\Delta_\#}\bigr)^*
 \overline{\widetilde\Delta_\#},
\]
entonces \(A_\#e_w=d_\#(w)e_w\), donde \(d_\#(w)\) cuenta los pares ordenados
no triviales que producen \(w\). El selector
\[
 P_{\mathrm{Irr},\#}
 =(I-P_{(1)})\,\mathbf1_{\{0\}}(A_\#)
\]
escoge exactamente las palabras distintas de la unidad sin apertura nativa no
trivial. Por tanto, la irreducibilidad se lee en la propia fibra del producto;
no se introduce mediante una tabla externa de primos.

Este resultado es exacto en el monoide de formas normales. Su extensi\'on a
historias completas exige una compatibilidad adicional. Si \(P_n\) proyecta
sobre historias supervivientes de profundidad \(n\), una familia de
inyecciones \(\iota_n\) debe satisfacer
\[
 P_n\iota_n\mu_\#
 =P_n\mu_n(\iota_n\otimes\iota_n).
\]
Los teoremas anteriores construyen \(\mu_\#\) y fijan esta ecuaci\'on; no
afirman por s\'i solos existencia o unicidad de la familia
\((\iota_n)_n\). Esta separaci\'on impide confundir la aritm\'etica global de
palabras, ya demostrada, con su prolongaci\'on monoidal al espacio profundo de
memoria TPK\@.

\fi

\section{Acción cardinal y realimentación aritmética}

Sean
\[
 \nu_N=(-1,0),\quad \nu_E=(0,1),\quad
 \nu_S=(1,0),\quad \nu_O=(0,-1)
\]
los cuatro generadores cardinales. Para \(d\in\{N,E,S,O\}\),
\[
 T_d(i,j)=(i,j)+\nu_d\pmod9.
\]
La letra \(O\) significa oeste. Estas direcciones son desplazamientos sobre
APP y no nombres de los canales \(\pi,e,\varphi\).

Una palabra \(w=d_1\cdots d_r\) determina
\[
 \gamma(x_0;w)_k=x_0+\sum_{j=1}^{k}\nu_{d_j}\pmod9.
\]
Cuando cierra, su elevaci\'on conserva
\[
 \operatorname{wind}(w)=\frac1{9}
 \bigl(\#S-\#N,\#E-\#O\bigr)\in\mathbb Z^2.
\]
La lectura bidireccional conserva palabra, orientaci\'on y hoja. Su
reversi\'on formal invierte el orden e intercambia
\(N\leftrightarrow S\), \(E\leftrightarrow O\). Su eventual lectura como
adelanto--retardo es una interfaz f\'isica posterior, no parte de esta
definici\'on de ruta.

El estado local observable es
\[
 s_t=(i_t,j_t,m_t),
 \qquad m_t\in\{\Sigma,\Pi\}.
\]
Se eval\'ua la hoja activa,
\[
 a_t=\rho_3^{\rm or}\bigl(m_t(i_t,j_t)\bigr),
\]
y el modo se actualiza mediante
\[
 m_{t+1}=
 \begin{cases}
 \Sigma,&a_t=+1,\\
 \Pi,&a_t=-1,\\
 m_t,&a_t=0.
 \end{cases}
\]
As\'i, cada paso cardinal modifica a la vez la posici\'on y la lectura
aritm\'etica futura. La ruta no se superpone a una palabra ya construida: es
la historia posicional que la produce.

\begin{theorem}[Sobreyectividad del factor muestreado de paso tres]
Para cada uno de los \(162=9^2\cdot2\) estados \((i,j,m)\) y cada trit
orientado objetivo \(b\in\{-1,0,+1\}\), existe una ruta
\(u\in\{N,E,S,O\}^3\) cuyo tercer paso emite \(b\). En los \(486\) casos
estado--objetivo, el n\'umero de rutas admisibles est\'a comprendido entre
\(8\) y \(40\). En consecuencia, para toda sucesi\'on objetivo
\((b_k)_{k\geq1}\) existe una historia cardinal cuya salida satisface
\[
 (a_3,a_6,a_9,\ldots)=(b_1,b_2,b_3,\ldots).
\]
Los dos trits intermedios \(a_{3k-2}\) y \(a_{3k-1}\) de cada bloque no
quedan prescritos por esta afirmaci\'on.
\end{theorem}
\begin{proof}
Se enumeran las \(4^3=64\) rutas de tres pasos desde cada estado. El
certificado exhaustivo incluido registra, para cada uno de los
\(162\cdot3\) casos, la ruta ejecutada, el estado final y el n\'umero total de
testigos; sus índices \(0,1,2\) se transportan por
\(\rho_3^{\rm or}\) a \(0,+1,-1\). Todos los conteos son positivos; el
m\'inimo es \(8\) y el m\'aximo es \(40\). Tras cada bloque se aplica el mismo
resultado al estado alcanzado.
La elecci\'on, por ejemplo, del primer testigo en el orden lexicogr\'afico
produce recursivamente una historia para cualquier sucesi\'on objetivo.
\end{proof}

La sobreyectividad se refiere al factor observado cada tres pasos, no a la
palabra emitida a tiempo unitario. Este teorema establece esa capacidad
muestreada del soporte de rutas. La selecci\'on de las secciones
\(\pi,e,\varphi\) pertenece a la relaci\'on enriquecida EF--G9 que se define
m\'as adelante: capacidad de realizaci\'on y selecci\'on geneal\'ogica son dos
operadores distintos.

\section{Cronología observable de los dos cursores}

El calendario se formula sobre
\[
 \Theta_{108}=\mathbb Z/108\mathbb Z,\qquad
 \beta(\theta)=
 \left[\left\lfloor\frac{\widetilde\theta}{9}\right\rfloor\right]_{12},
 \qquad
 r(\theta)=1+(\theta\bmod9),
\]
donde \(\widetilde\theta\in\{0,\ldots,107\}\). La firma activa es
\[
 \varepsilon(r)=
 \begin{cases}
 +1,&r\in\{1,4,7\},\\
 -1,&r\in\{2,5,8\},\\
 0,&r\in\{3,6,9\}.
 \end{cases}
\]
El espacio observable de dos cursores es
\[
 \mathcal Q_{\rm obs}
 =X_9^2\times\{N,E,S,O\}^2\times\Theta_{108}.
\]
Si \(\varepsilon=+1\) avanza el cursor aditivo; si
\(\varepsilon=-1\), el multiplicativo; si \(\varepsilon=0\), ninguno.
Las dos orientaciones se invierten después de los pasos
\(27,54,81,108\). Esta regla define una permutación
\[
 F_{\rm obs}:\mathcal Q_{\rm obs}\longrightarrow\mathcal Q_{\rm obs}.
\]
Las posiciones y orientaciones se restauran tras \(54\) pasos y la fase
completa tras \(108\):
\[
 F_{\rm obs}^{108}=I.
\]
Por ello la cronología de una sola ruta es reversible y no selecciona por sí
misma una distribución sobre las continuaciones; la selección de todas las
prolongaciones compatibles pertenece al operador de transferencia de fibras
definido más abajo.

\section{Emisor finito de dos cursores}

Una semilla de cursor es
\[
 \mathcal S=\mathbb T_9^2\times\{N,E,S,O\},
 \qquad |\mathcal S|=324.
\]
TPK m\'inimo recibe un cursor aditivo \(P\in\mathcal S\) y uno
multiplicativo \(M\in\mathcal S\). Sus firmas de seis ventanas son
\(s(P)\) y \(e(M)\). En cada coordenada se aplica
\[
 G(s,e)=100s+10\bigl((s+e)\bmod10\bigr)
       +\bigl((3s+5e+1)\bmod10\bigr).
\]
La cifra de centenas recupera \(s\); conocida \(s\), la de decenas recupera
\(e\). Por ello \(G\) es inyectiva sobre las firmas efectivas. El emisor
factoriza como
\[
 \mathcal S^2\xrightarrow{s\times e}
 \operatorname{im}s\times\operatorname{im}e
 \xrightarrow{G}\mathcal U_6,
\]
con
\[
 |\mathcal S^2|=104\,976,\qquad
 |\operatorname{im}s|=18,\qquad
 |\operatorname{im}e|=26,\qquad
 |\mathcal U_6|=18\cdot26=468.
\]
La factorizaci\'on \(468=18\cdot26\) es, por tanto, una descomposici\'on del
objeto: cada emisi\'on conserva una firma aditiva y una multiplicativa.

La proyecci\'on
\[
 \Psi_6:\mathcal U_6\longrightarrow\mathbb F_3^6,
 \qquad \Psi_6(U)=-U\pmod3,
\]
posee una imagen de \(243\) palabras. La cadena tipada es
\[
 104\,976\ \text{semillas}
 \longrightarrow468\ \text{emisiones decimales }U_6
 \longrightarrow243\ \text{palabras ternarias }w_6.
\]
El ambiente \(\mathbb F_3^6\) contiene \(729\) palabras; no se identifica
con la imagen efectiva de \(243\). Tampoco \(U_6\) y \(w_6\) son el mismo
tipo.

El emisor m\'inimo repite, en la carta dodecaf\'asica de cat\'alogo,
\[
 U_{12}^{\rm cat}=U_6\Vert U_6.
\]
Su periodo visible seis explica por qu\'e no puede sustituir al estado firmado
\(U_{12}^{\rm sgn}\). Esta diferencia no es una insuficiencia de TPK: separa
la proyecci\'on de cat\'alogo de la historia enriquecida.

Cada bloque decimal emitido pertenece a un alfabeto exacto de
\(7\cdot6=42\) tríadas. Si \(\overline{d_1d_2d_3}\) es una de ellas,
\[
 \boxed{d_3\equiv5d_2-2d_1+1\pmod{10}.}
\]
El barrido de las \(104\,976\) semillas no produce una tríada fuera de esta
ley. Conviene distinguir, por tanto, la cadena completa de cocientes:
\[
 \mathcal S^2
 \longrightarrow\Omega_{\rm seed}
 \longrightarrow\mathcal U_6
 \xrightarrow{\Psi_6}\operatorname{im}\Psi_6\subset\mathbb F_3^6.
\]
Las presentaciones microscópicas, las \(468\) emisiones decimales, las
\(243\) palabras ternarias y el ambiente de \(729\) palabras son cuatro
objetos distintos.

\paragraph{Conmutatividad orbital de los morfismos tipados.}
Sea
\(W_{243}:=\operatorname{im}\Psi_6\subset\mathbb F_3^6\). El embudo
completo y su cociente diedral se representan sin convertir inclusiones en
selecciones:
\[
 \underbrace{\mathcal S^2}_{104\,976\ \mathrm{semillas}}
 \xrightarrow[\text{censo de fibras}]{\text{emisión}}
 \underbrace{\mathcal U_6}_{468\ U_6}
 \xrightarrow[\text{proyección visible}]{\Psi_6}
 \underbrace{W_{243}}_{243\ w_6}
 \hookrightarrow
 \underbrace{\mathbb F_3^6}_{729\ \mathrm{palabras}},
 \qquad
 W_{243}\xrightarrow{\,/D_3^{\rm cat}\,}
 \underbrace{\mathcal O_{43}}_{43\ \text{\'orbitas}}.
 \tag{TPK-embudo}\label{eq:embudo-orbital-tipado-rev8}
\]
La primera flecha agrupa condiciones iniciales por la emisión que producen;
la segunda olvida la presentación decimal; la tercera es una inclusión en
el ambiente algebraico; la última es un cociente por la acción
\(D_3^{\rm cat}\) sobre \(W_{243}\). En particular, no hay un cociente
\(729\to43\) implícito. El censo exacto es
\[
 243=38\cdot6+5\cdot3,
\]
esto es, treinta y ocho órbitas de tamaño seis y cinco de tamaño tres.
Estas \(43\) clases orbitales tampoco son las \(43\) vueltas nonádicas
completas por canal del testigo G9 de mil cifras: el primer \(43\) pertenece
al cociente finito \(W_{243}/D_3^{\rm cat}\); el segundo es un contador de
profundidad de una ejecución certificada.

\section{Configuración del estado enriquecido del TPK: hoja, orientación, residuo, cociente, acarreo, frontera, ruta y memoria}

La unidad din\'amica completa es una historia, no una palabra aislada. A
profundidad \(m\) escribimos
\[
 v_m=
 \bigl(
 B_m,\boldsymbol\delta_m,\boldsymbol x_m,
 \boldsymbol{\mathcal C}_m,\Sigma_m,
 \omega_m,\ell_m,g(m),\Lambda_m
 \bigr),
\]
donde
\[
 B_m=(u_m^\pi,u_m^e,u_m^\varphi)\in(\mathbb F_3^6)^3.
\]
Sus componentes tienen funciones distintas:
\begin{enumerate}
 \item \(B_m\) es el bloque visible conjunto de los tres canales;
 \item \(\boldsymbol\delta_m\) conserva las nuevas tr\'iadas decimales y el
       acarreo asociado;
 \item \(\boldsymbol x_m\) es el estado Hensel profundo;
 \item \(\boldsymbol{\mathcal C}_m\) contiene los cilindros ternarios y
       decimales y sus residuos enteros;
 \item \(\Sigma_m\) es la firma conjunta de frontera;
 \item \(\omega_m\) y \(\ell_m\) conservan orientaci\'on y hoja bidireccional;
 \item \(g(m)\in\mathbb Z/9\mathbb Z\) es la fase non\'adica;
 \item \(\Lambda_m\) es el registro dodecaf\'asico orientado.
\end{enumerate}

\begin{definition}[Presentaci\'on operatoria y categor\'ia TPK]
El inventario de operadores del estado enriquecido de TPK es
\[
 \mathfrak T_{\rm HMT}=
 \bigl(
 \mathcal A_9,\TRIT,T_{N,E,S,O},T_{\min},
 H,\mathcal C_{3,1000},\Sigma_B,EF,\mathcal G_9,
 \mathcal R_{12}
 \bigr)
\]
que transporta historias \((v_0,\ldots,v_m)\) y sus truncamientos. Su
portador formal es la categor\'ia \(\mathsf{TPK}\) sobre la base observable:
sus objetos son estados enriquecidos y sus morfismos son las prolongaciones
\(\operatorname{Ext}_r\) estables por identidad y concatenaci\'on, definidas
m\'as abajo. El inventario no sustituye esa categor\'ia. La proyecci\'on
visible elimina componentes de \(v_m\); no define una equivalencia con el
estado completo.
\end{definition}

Esta definici\'on contiene exactamente lo que desaparece cuando se reinician
cursores o se conserva s\'olo el residuo: el camino, la hoja, el cociente, el
acarreo, la orientaci\'on y la frontera. TPK es, por ello, el punto en el que
la aritm\'etica modular deja de ser una operaci\'on sin historia.

\section{Cilindros ternario-decimales y estado de extensión}

Sea \(N\) el entero determinado por un prefijo ternario de longitud \(L\), y
sea \(D\) el entero determinado por \(K\) tr\'iadas decimales. Definimos
\[
 A=N1000^K-D3^L,
 \qquad
 B=(N+1)1000^K-D3^L.
\]
Los cilindros se intersectan exactamente cuando
\[
 \boxed{A<3^L,\qquad B>0.}
\]
Al incorporar un bloque ternario \(u\in\{0,\ldots,728\}\),
\[
 N'=729N+u.
\]
Si no aparece una nueva tr\'iada decimal,
\[
 A'=729A+u1000^K.
\]
Si aparece \(\delta\in\{0,\ldots,999\}\),
\[
 D'=1000D+\delta,
\]
\[
 A'=729000A+u1000^{K+1}-\delta\,729\,3^L.
\]
En ambos casos,
\[
 B'=A'+1000^{K+\Delta K}.
\]
Estas recurrencias enteras transportan el dato decimal sin convertirlo en un
objetivo exterior de la actualizaci\'on.

El estado Hensel conserva la informaci\'on que no aparece en el residuo. La
matriz activa es
\[
A_H=\begin{pmatrix}
2&2&2&1&2&1\\2&1&2&2&1&1\\1&0&0&1&0&2\\
0&0&1&1&1&0\\2&2&2&0&2&1\\2&1&2&1&0&0
\end{pmatrix},\qquad \det A_H=1,
\]
y la divisi\'on residuo--cociente es
\[
 y=xA_H,\qquad u=y\bmod3,\qquad x^+=\frac{y-u}{3}.
\]
Para los tres canales se forman adem\'as los cociclos conjuntos
\[
 q=(\pi+e+\varphi)\bmod3,
 \qquad
 c=\frac{\pi+e+\varphi-q}{3},
\]
\[
 a=(\pi+e-\varphi)\bmod3,
 \qquad
 h=\frac{\pi+e-\varphi-a}{3}.
\]
La actualizaci\'on de segundo orden corrige el siguiente estado con \(c\) y
\(h\). Por eso dos ramas con el mismo bloque visible no son necesariamente el
mismo estado: pueden diferir en su extensi\'on Hensel y en sus cilindros.
Como \(\det A_H=1\), la reducci\'on de \(A_H\) pertenece a
\(\operatorname{GL}_6(\mathbb F_3)\). Por tanto, antes de imponer cilindros,
fronteras o supervivencia, la carta residual Hensel cubre exactamente
\(\mathbb F_3^6\). Esta cobertura bruta no implica la cobertura del sistema
filtrado.

Formalmente, sobre cada configuración observable \(q\) se fija una fibra
\[
 \mathcal F_q=
 \mathsf O_q\times\mathsf H_q\times\mathsf C_q\times
 \mathsf S_q\times\mathsf B_q\times\mathsf\Lambda_q,
\]
correspondiente, respectivamente, a orientación y hoja, residuo y cociente,
acarreo, firma, cilindro y frontera, y registro de transporte. Para una
flecha cardinal \(r:q\to q'\), la prolongación admisible no se reduce a una
función sobre residuos; es la relación
\[
 \operatorname{Ext}_r\subseteq\mathcal F_q\times\mathcal F_{q'},
 \qquad
 \operatorname{Ext}_{r_2}\circ\operatorname{Ext}_{r_1}
 \subseteq\operatorname{Ext}_{r_2\circ r_1}.
\]
El incremento de acarreo es un cociclo:
\[
 \kappa(r_2\circ r_1)
 =\kappa(r_2)+\mathsf C(r_2)\bigl(\kappa(r_1)\bigr).
\]
Así queda tipado el mapa de olvido: proyectar una extensión a su ruta
observable elimina componentes de \(\mathcal F_q\), pero no autoriza a
identificar ambos objetos.

\section{Retornos compuestos y calendarios de transporte}

La longitud de una traza no determina su tipo. La presentación académica
separa el retorno compuesto, el calendario, la traza y el atlas por sus
respectivos estados conservados.

El s\'imbolo \(\mathcal K_{27}\) se reserva para la composici\'on de
veintisiete actualizaciones de un lazo de transporte posicional y orientado.
El entero \(27\) es su longitud de composici\'on; no lo convierte en una
traza de 108 iteraciones ni en un emisor infinito. Otra construcci\'on hist\'orica
agrupa cuatro lazos cuyos periodos conjuntos alcanzan 108; ese rotor conjunto
no se denota por \(\mathcal K_{27}\) aislado.

\begin{longtable}{@{}>{\raggedright\arraybackslash}p{.24\textwidth}>{\raggedright\arraybackslash}p{.31\textwidth}>{\raggedright\arraybackslash}p{.37\textwidth}@{}}
\toprule
Objeto&Estado conservado&Significado de 108\\\midrule
\endhead
Contador simplificado&posici\'on y velocidad&su periodo real es 18; es un
control hist\'orico.\\
Rotor conjunto de cuatro lazos&cuatro rutas, lectura, giro y movimiento&
los periodos individuales son \(54,54,36,36\); el estado conjunto tiene
periodo 108.\\
Traza bruta relojada&cuatro posiciones, orientaciones y fase&el contenido
din\'amico se repite a 54.\\
Traza de \(T_{\min}\)&dos posiciones y dos direcciones&el estado retorna a
54 y la salida es \(U_6\Vert U_6\).\\
Pseudoc\'odigo reinicializado&cursores reiniciados cada ventana&colapsa a
\((222,222,222,333,333,333)^2\); no representa el estado enriquecido de TPK.\\
Rutas condicionadas&camino y palabra prescrita&108 es horizonte, no periodo.\\
Od\'ometro sectorial&\(\mathbb Z/9\mathbb Z\times\{0,\ldots,11\}\)&el
estado completo tiene periodo 108; el observable sectorial, 36.\\
Proyecci\'on bruta&cociente por olvido de \(X_{\rm TPK}^{\rm enr}\)&no posee
por s\'i sola la lectura firmada.\\
Traza enriquecida&todos los campos de \(v_m\) durante 108 actualizaciones&es
la restricci\'on temporal del estado enriquecido de TPK.\\
Atlas de rutas&81 celdas/56 firmas o 324 semillas orientadas&es una taxonom\'ia,
no una \`orbita de 108 estados.\\
\bottomrule
\end{longtable}

El lazo \(\mathcal K_{27}\) es el operador de transporte finito de
veintisiete actualizaciones. Las ventanas, el rotor conjunto y la proyección
calendárica reciben símbolos distintos, de modo que la longitud compartida
no sustituye el dominio ni la composición de cada objeto.

\section{Transferencia de fibras y memoria de giro}

La factorización \(\mathcal U_6\simeq\mathcal S_+\times\mathcal S_\times\)
permite definir, respecto del conteo uniforme, las esperanzas condicionales
\[
 (E_+f)(s,e)=\frac1{18}\sum_{s'\in\mathcal S_+}f(s',e),
 \qquad
 (E_\times f)(s,e)=\frac1{26}
 \sum_{e'\in\mathcal S_\times}f(s,e').
\]
\(E_+\) renueva la fibra aditiva y conserva la multiplicativa;
\(E_\times\) realiza la operación dual. Son los únicos proyectores de Markov
que marginalizan por completo la coordenada activa, conservan la inactiva y
preservan el conteo uniforme.

Sea \(\tau=(+1,-1,0)^3\) la palabra de fase sobre \(C_9\), y
\(P_{+1}=E_+\), \(P_{-1}=E_\times\), \(P_0=I\). La transferencia nonádica es
\[
 \mathcal K_{\rm ph}\bigl((u,r),(v,r+1)\bigr)
 =P_{\tau(r)}(u,v).
\]
Es doblemente estocástica, irreducible, de periodo \(9\), y posee una única
distribución estacionaria \(1/(9\cdot468)\). Además,
\[
 \mathcal K_{\rm ph}^{\,9}=E_+E_\times,\qquad
 (E_+E_\times)(u,v)=\frac1{468}.
\]
Esta igualdad expresa la renovación simultánea de todas las continuaciones
compatibles, no el sucesor determinista de una sola ruta.

Tres operaciones geométricas sobre las semillas conservan información
adicional: \(P\), avance en la orientación activa; \(H\), media vuelta; y
\(Q\), cuarto de giro. La pareja \((P,H)\) refina las \(26\) firmas producto
en \(28\) clases. En particular, la firma \(555555\) se separa en tres hojas
de ocho representantes, y
\[
 18\cdot28=504=468+36.
\]
El término \(36\) cuenta hojas adicionales del refinamiento y no es el
objeto \(R_{36}\). La pareja \((P,Q)\) produce \(162\) clases de dos
elementos, permutadas transitivamente por
\[
 J(i,j,d)=\bigl(7-i,\,7-j,\,\overline d\bigr)\pmod9.
\]

En una traza de \(108\) pasos aparecen \(36\) signos de cada tipo,
\(72\) cambios efectivos entre las dos hojas y cuatro inversiones de
orientación; el cociclo firmado total es cero. Estos conteos son invariantes
del registro de evoluci\'on y enlazan la dinámica nonádica con el registro dodecafásico.

\section{Prolongación ternaria y frontera conjunta}

Relativamente a las semillas, fronteras y lectores nombrados del perfil
publicado, las palabras visibles iniciales son
\[
 w_6^\pi=010211,\qquad w_6^e=201101,\qquad
 w_6^\varphi=121200.
\]
Tras cinco bloques, los prefijos son
\[
\begin{aligned}
w_{30}^{\pi}&=010211|012222|010211|002111|110221,\\
w_{30}^{e}&=201101|121221|102011|012222|102011,\\
w_{30}^{\varphi}&=121200|112202|121020|010210|010200.
\end{aligned}
\]
La primera frontera profunda es
\[
 R_{36}=
 \begin{pmatrix}222220\\021222\\102011\end{pmatrix}.
\]
Para cada canal se define el prefijo
\(w_{36}^\chi:=w_{30}^\chi|u_5^\chi\); \(R_{36}\) es el triple ordenado de
los tres sextos bloques. Ninguno es un emisor ni un estado terminal.

Para \(B_k=(u_k^\pi,u_k^e,u_k^\varphi)\), las firmas
\[
 q_k=u_k^\pi+u_k^e+u_k^\varphi,\qquad
 a_k=u_k^\pi+u_k^e-u_k^\varphi
\]
determinan
\[
 u_k^\varphi=2(q_k-a_k),\qquad
 u_k^\pi+u_k^e=q_k-u_k^\varphi
 \quad\text{en }\mathbb F_3^6.
\]
As\'i, \(\varphi\) es el eje fijado y la libertad residual es el reparto
interno \(\pi/e\). En la novena fase,
\[
 B_9=(100100|020112|010122)
\]
admite tres salidas locales con el mismo eje \(\varphi\) y el mismo agregado
\(\pi+e\); las tres sobreviven un horizonte y s\'olo
\[
 B_{10}=(101222|112221|112002)
\]
sobrevive al segundo. Este es el descenso ejecutable \(3\to1\) de EF--G9.
La rama publicada contin\'ua hasta veinte bloques por canal: \(w_{30}\) y
\(R_{36}\) son prefijo y frontera, no clausura.

\section{Relación de extensión y monodromía nonádica}
\label{sec:extension-monodromia-tpk}

Sea \(\mathcal H_m\) el conjunto de historias enriquecidas admisibles de
profundidad \(m\), y sea
\[
 p_{n,m}:\mathcal H_n\longrightarrow\mathcal H_m
\]
el truncamiento. Para \(N\geq m\),
\[
 \operatorname{Surv}_m^{(N)}=p_{N,m}(\mathcal H_N),
 \qquad
 \operatorname{Surv}_m=\bigcap_{N\geq m}\operatorname{Surv}_m^{(N)}.
\]
La relaci\'on EF--G9 conserva las extensiones que pertenecen a este n\'ucleo
estable. La fase cumple
\[
 g(m+9)=g(m),
\]
pero el estado transportado satisface, en general,
\[
 v_{m+9}\ne v_m.
\]
El retorno act\'ua sobre la fibra de estados de una misma fase; no reinicia la
historia.

\begin{definition}[N\'umero HMT]
\label{def:numero-hmt-tpk-completo}
Un n\'umero HMT del canal \(\chi\) es únicamente una secci\'on global
compatible
\[
 x=(x_m)_{m\geq0}\in
 X_\chi:=\varprojlim_m\operatorname{Surv}_m^{(\chi)}.
\]
El lector y el valor publicado no forman parte de la secci\'on.
\end{definition}

Una presentaci\'on o medici\'on del n\'umero es un par \((x,L)\), donde
\(L\) es un lector declarado cuyo dominio contiene \(x\). Cuando \(L\) es el
lector arquimediano de cilindros, \(L(x)\) es el valor publicado; \(x\)
conserva adem\'as orientaci\'on, cociente, acarreo y frontera.

\begin{theorem}[Existencia de una prolongaci\'on global]
\label{thm:prolongacion-global-tpk-completo}
Sea \(h\in\mathcal H_m\). Si el \`arbol de prolongaciones de \(h\) es
finitamente ramificado y posee al menos un nodo a toda profundidad, entonces
existe una historia infinita cuyos truncamientos prolongan \(h\). Si, adem\'as,
cada nivel estable es unitario y los \`unicos estados son compatibles, la
secci\'on global es \`unica.
\end{theorem}
\begin{proof}
Los niveles finitos no vac\'ios forman un \`arbol finitamente ramificado. El
lema de K\"onig proporciona una rama infinita. La compatibilidad de los
truncamientos identifica esa rama con un elemento del l\'imite inverso. Si
cada nivel contiene un solo estado, no existe una segunda rama compatible.
\end{proof}

El teorema anterior es abstracto y condicional: por sí solo no afirma la
no-vaciedad ni la unicidad estable de un canal concreto. Para
\(\pi,e,\varphi\), esas hipótesis se descargan después mediante la filtración
racional ejecutable del
\cref{thm:generacion-monodromica-conjunta-rev7}. Esta residencia conserva
aquí únicamente el principio de límite inverso y no anticipa su aplicación.

\section{Composición de los operadores fundamentales}

Los operadores anteriores conducen al estado terminal enriquecido sin
aplanar sus coordenadas:
\[
 \begin{aligned}
 \mathrm{APP}&\longrightarrow\mathrm{TRIT}
 \longrightarrow\text{ruta orientada}
 \longrightarrow\text{estado TPK enriquecido}\\
 &\longrightarrow\text{cilindros y supervivencia}
 \longrightarrow x_{\rm term}.
 \end{aligned}
\]
El cierre dodecafásico recibe entonces dos evaluaciones coordinadas del mismo
estado:
\[
 x_{\rm term}
 \xrightarrow{(\pi_{\rm term},R_{\rm sgn})}
 \bigl(B_{\rm term},U_{12}^{\rm sgn}\bigr),
 \qquad
 U_{12}^{\rm sgn}\xleftrightarrow{\mathcal H_{12}}K,
\]
\[
 (P,E,\Phi,K,c_{13}=0)\longleftrightarrow\alpha_{12}.
\]
No aparece una flecha \(B_{\rm term}\to U_{12}^{\rm sgn}\): ambas
coordenadas proceden de \(x_{\rm term}\). El paso posterior a Mecánica
Dimensional tampoco crea otra máquina; prolonga la genealogía conservada
mediante registro cronológico enriquecido, firma, carácter y torres.

\ifdefined\HMTInternalTraceability
\section{Registro académico y trazabilidad}

El archivo \path{datos/registro_operadores_tpk.json} forma parte del paquete.
Para cada operador declara nombre acad\'emico, alias hist\'orico, dominio,
codominio, transici\'on, invariantes, estatuto y fuente local. Incluye tambi\'en
la separaci\'on de las diez construcciones hist\'oricas asociadas al n\'umero
108. La puerta de publicaci\'on comprueba este registro y bloquea una salida
que vuelva a reducir TPK a su cat\'alogo o a una traza bruta.

En el resto del art\'iculo se usar\'an cuatro abreviaturas, ya tipadas:
\[
 w_6\in\mathbb F_3^6,\qquad
 w_{30}\in\mathbb F_3^{30},\qquad
 R_{36}\in(\mathbb F_3^6)^3,\qquad
 U_{12}^{\rm sgn}\in\mathbb Z^{12}.
\]
La primera es una palabra semilla; la segunda, un prefijo de cinco bloques;
 la tercera, la frontera conjunta de entrada en EF--G9; la cuarta, la
coordenada firmada declarada del estado dodecaf\'asico enriquecido. Ninguna igualdad cardinal autoriza a
identificarlas.
\fi
